\documentclass[10pt,a4paper]{amsart}
\usepackage[margin=19mm]{geometry}
\usepackage{amsmath,amssymb,mathtools}
\usepackage[scr=rsfs,cal=euler]{mathalfa}
\usepackage{xfrac}
\usepackage[unicode,hidelinks,bookmarks=false]{hyperref}
\newcommand{\ie}{i.e.\ }
\newcommand{\cf}{cf.\ }
\newcommand{\resp}{respectively\ }
\newcommand{\Subsection}{\subsection}
\providecommand{\nobreakdash}{\nobreak\hbox{-}}
\setlength{\emergencystretch}{2em}
\multlinegap0pt
\usepackage[shortlabels]{enumitem}
\usepackage{amssymb}
\usepackage{xcolor}
\usepackage{float}
\usepackage{tikz}
\newtheorem{lem}{Lemma}
\newtheorem{cor}{Corollary}
\newtheorem{thm}{Theorem}
\DeclareMathOperator{\cof}{cof}
\title{The Cofinality of Generating Families}
\author{Paul Gartside, Thomas Gilton}
\date{Source: arXiv:2602.02211v2 (CC BY 4.0)}
\begin{document}
\maketitle

\noindent\emph{Source and licence.} The Cofinality of Generating Families. Version arXiv:2602.02211v2 (CC BY 4.0), distributed by arXiv under the Creative Commons Attribution 4.0 International licence, http://creativecommons.org/licenses/by/4.0/, as recorded in the arXiv licence metadata for this exact version and shown on the abstract page https://arxiv.org/abs/2602.02211v2. Full licence notice: \texttt{licenses/arxiv-2602.02211-CC-BY-4.0.txt}.\par\smallskip

\noindent\textit{Editorial scope.} Counted unit: the article's thm th:K\_many (source0.tex lines 863--875). The article's complete original proof of it is delivered with it (source0.tex lines 900--930, one \texttt{proof} environment of 31 lines, which the article places away from the statement). No mathematics is written, repaired or re-derived: the statement and the proof are the article's own lines, in the article's own notation, and no retained line is substituted or re-flowed.  Retained uncounted as the source-local context the proof needs: lines 227--238; lines 374--377; lines 429--431; lines 774--784; lines 831--835.  Nothing was omitted from any retained span, so the package carries no omission record.  No retained line contains a citation, so the wrapper carries no bibliography and no bibliography entry.  No retained line contains a figure environment, an image-inclusion command or drawing code, so no image is delivered and the package declares no asset dependency.  Editorial formatting only, in addition: an AMS \texttt{amsart} preamble that restates the article's own declarations and macros used by the retained lines, namely the environments \texttt{lem}, \texttt{cor}, \texttt{thm} on separate counters, together with the standard packages those lines need.  The retained environments print in this document as Lemma 1, Corollary 1, Theorem 1 in order of appearance, whereas the article prints the same items with the numbering of its own full document.  The complete native source of the article as distributed in the arXiv e-print for this exact version is bundled unchanged as \texttt{sources/193834/source0.tex}.
\begin{lem} \label{l:sam1} 
Let $\kappa$ be an uncountable cardinal. Then
\begin{enumerate}[noitemsep,topsep=1pt]
\item[(i)]  $\kappa \le \mathop{cov}(\kappa) \le \kappa^\omega$, and $\mathop{cov}(\kappa^\omega)=\kappa^\omega$, 

\item[(ii)] if $\cof(\kappa)=\omega$ then $\mathop{cov}(\kappa) > \kappa$,

\item[(iii)] if $\kappa=\mu^+$ then     $\mathop{cov}(\kappa) = \mathop{cov}(\mu)\cdot\kappa$,  and

\item[(iv)] if $\kappa$ is a limit and $\cof(\kappa)>\omega$  then $\mathop{cov}(\kappa)=\lim \{\mathop{cov}(\mu) : \mu<\kappa\}$. 
\end{enumerate}
\end{lem}

\begin{lem}\label{lem:cccPreservesCapturing}
    Let $\kappa$ be an uncountable cardinal.
        Let $\mu=\mathop{cov}(\kappa)$ and let $\mathbb{P}$ be c.c.c. Let $\dot{X}$ be the $\mathbb{P}$-name for the set of countable subsets of $\kappa$ in the extension. Then $\mathbb{P}$ forces that $\mu=\cof(\dot{X},\subseteq)$.
\end{lem}

\begin{cor}\label{cor:ScalesIncreaseSam}
    Suppose that $\kappa$ is a singular cardinal of countable cofinality. If $\kappa$ admits a scale of regular length $\lambda$ then $\mathop{cov}(\kappa)\geq\lambda$. 
\end{cor}

\begin{thm}\label{th:k=sam}
Let $M$ be a non $\sigma$-compact, separable metrizable space.
\begin{enumerate}[noitemsep,topsep=1pt]
\item[(i)] If $M$ is totally imperfect but not locally small then 
$k(M)=\mathop{cov}(kc(M))\cdot\mathfrak{b}$.

\item[(ii)] If $M$ is locally small then 
$
k(M)=\lim_{\mu < kc(M)} \mathop{cov}(\mu)\cdot\mathfrak{b}$. 
\end{enumerate}
\end{thm}

\begin{thm}\label{th:powers} 

    Let $N$ be separable metrizable but not compact, and set $M=N^\omega$. Then $kc(M) \ge \mathfrak{d}$ and $k(M)=kc(M)$.     
    Further, if $|N|<\mathfrak{c}$ then $k(M)=kc(M)=\mathop{cov}(|N|)\cdot \mathfrak{d}$. 
\end{thm}

\begin{thm}\label{th:K_many}
Fix an integer $K\ge 2$. 
Let $\lambda=\aleph_{\omega\cdot K}$. 

Modulo large cardinals, it is consistent that $\mathfrak{b}=\aleph_1$, 
$\mathop{cov}(\lambda^+)=\lambda^{+K}$, and  
for $n=1,\ldots,K$ there are separable metrizable spaces, $M_n$, such that $kc(M_n)=\lambda^+$ while $k(M_n)=\lambda^{+n}$.

Hence in this model, $\kappa:=\lambda^+ \le \mathfrak{c}$ is uncountable and
\[\{k(M):kc(M)=\kappa\}=
[\lambda^+,\lambda^{+K}]=
[\kappa\cdot\mathfrak{b},\mathop{cov}(\kappa)\cdot\mathfrak{b}], \ \text{  so }\sigma_k(\kappa)=K.\] 
\end{thm}

\begin{proof}[Proof of  Theorem~\ref{th:K_many}]
Fix $K\ge 2$. 
Start with the model of Gilton and Cummings above. For each $1\leq n\leq K$, from (2), (3) together with Corollary~\ref{cor:ScalesIncreaseSam}   conclude that 
 $\mathop{cov}(\aleph_{\omega\cdot n})$ is equal to $\aleph_{\omega\cdot K+n}$. 
    Now do random real forcing to increase the value of $\mathfrak{c}$ strictly above $\aleph_{\omega\cdot K+K}$, while keeping $\mathfrak{b}$ and $\mathfrak{d}$ fixed at $\omega_1$. By Lemma~\ref{lem:cccPreservesCapturing}, this forcing doesn't change the relevant values of the covering numbers. 
Thus, in the resulting model, (a) $\mathfrak{b}=\mathfrak{d}=\omega_1$;
    (b) $\mathfrak{c}>\aleph_{\omega\cdot K+K}$; and
    (c) for each $1\leq n\leq K$, $\mathop{cov}(\aleph_{\omega\cdot n})=\aleph_{\omega\cdot K+n}$.

Let $\lambda=\aleph_{\omega\cdot K}$. 
Let $N$ be a subset of the reals of size $\aleph_\omega$, and set $M_\infty=N^\omega$. For each $1\le n \le K$, let $N_n$ have size $\aleph_{\omega\cdot n}$  but not be locally small, and set $M_n = M_\infty \oplus N_n$. 

We verify the claims of Theorem~\ref{th:K_many}. Item (a) simply asserts that $\mathfrak{b}=\omega_1$.
From Lemma~\ref{l:sam1}(iii) and from item (c), we see 
\[
\mathop{cov}(\lambda^+)=\mathop{cov}(\aleph_{\omega\cdot K+1})=\mathop{cov}(\aleph_{\omega\cdot K})\cdot \aleph_{\omega\cdot K+1}=\aleph_{\omega\cdot K+K}=\lambda^{+K}.
\]


The compact-covering number of $M_n$, $kc(M_n)$, is the larger of $kc(M_\infty)$ and $kc(N_n)$. 
From Theorem~\ref{th:powers} we know $kc(M_\infty)=\mathop{cov}(|N|)\cdot\mathfrak{d}=\mathop{cov}(\aleph_\omega)\cdot\aleph_1=\aleph_{\omega\cdot K+1}$. 
And since $|N_n|=\aleph_{\omega\cdot n} < \mathfrak{c}$ (by (b)), we have that $N_n$ is totally imperfect (since it has no compact sets of size $\mathfrak{c}$, on account of having size $<\mathfrak{c}$). Hence, 
$kc(N_n)=|N_n|=\aleph_{\omega\cdot n}$. 
Therefore $kc(M_n)=\aleph_{\omega\cdot K+1}=\lambda^+$, as required.

The $k$-ness number of $M_n$, $k(M_n)$, is the larger of $k(M_\infty)$ and $k(N_n)$. 
From Theorem~\ref{th:powers} we know 
$k(M_\infty)=kc(M_\infty)=\aleph_{\omega\cdot K+1}$. 
While from Theorem~\ref{th:k=sam}, as $kc(N_n)=\aleph_{\omega\cdot n}<\mathfrak{c}$ (by (b)), we have $k(N_n)=\mathop{cov}(\aleph_{\omega\cdot n})=\aleph_{\omega\cdot K+n}$. 
Hence $k(M_n)=\aleph_{\omega\cdot K+n}=\lambda^{+n}$, as required.    
\end{proof}

\end{document}
